\section{最適化・意思決定}

\subsection{平成24年度 問13}
\label{q:h24-reg-apt-13}

\noindent\textbf{原典：} sources/H24/aptitude.pdf、PDF物理11頁。

\textbf{公式正答：} 3（single）。

\par\textit{原典本文の校正と独自解説は未作成。分類のみの登録項目。}

\subsection{令和2年度 問15}
\label{q:r02-reg-apt-15}

\noindent\textbf{原典：} sources/R02/aptitude.pdf、PDF物理15頁。

\textbf{公式正答：} 1（single）。

\par\textit{原典本文の校正と独自解説は未作成。分類のみの登録項目。}

\section{品質管理・統計}

\subsection{平成29年度 問7}
\label{q:h29-reg-apt-07}

\noindent\textbf{原典：} sources/H29/aptitude.pdf、PDF物理7頁。

\textbf{公式正答：} 4（single）。

\par\textit{原典本文の校正と独自解説は未作成。分類のみの登録項目。}

\section{暗号・認証・攻撃対策}

\subsection{令和元年度 問9}
\label{q:r01-reg-apt-09}

\noindent\textbf{原典：} sources/R01/aptitude.pdf、PDF物理10頁。

\textbf{公式正答：} 2（single）。

\par\textit{原典本文の校正と独自解説は未作成。分類のみの登録項目。}

\section{アルゴリズム・再帰}

\subsection{令和元年度再試験 問3}
\label{q:r01-ret-apt-03}

\noindent\textbf{原典：} sources/R01\_retest/aptitude.pdf、PDF物理3頁。

\textbf{公式正答：} 1（single）。

\par\textit{原典本文の校正と独自解説は未作成。分類のみの登録項目。}

\section{公益・説明責任}

\subsection{平成26年度 問6}
\label{q:h26-reg-apt-06}

\noindent\textbf{原典：} sources/H26/aptitude.pdf、PDF物理6頁。

\textbf{公式正答：} 5（single）。

\par\textit{原典本文の校正と独自解説は未作成。分類のみの登録項目。}

\section{社会的責任・安全の標準}

\subsection{平成23年度 問11}
\label{q:h23-reg-apt-11}

\noindent\textbf{原典：} sources/H23/aptitude.pdf、PDF物理12頁。

\textbf{公式正答：} 3（single）。

\par\textbf{検討未確定：} Astra Review 1後の校正・解説例。解説はAstra査読待ち。①政府除外、④原則列挙、②⑤歴史的経緯を要再確認。要旨のみ校正済みで全文の校正完了ではない。

\paragraph{問題の要旨} 社会的責任の手引きと認証規格の区別を問う。公式正答③を保持し、他の記述の解釈上の疑義を別に残す。

\Needspace{4\baselineskip}\paragraph{確認した条件・記述（独自要約）}\mbox{}\par

\noindent\textbf{standard：} ISO26000\par

\noindent\textbf{exam\_reference\_year：} 2011\par

\paragraph{問われる核心} 技術者倫理・意思決定

\paragraph{必要知識} 社会的責任、要求事項とガイダンス、認証の適用範囲

\paragraph{考え方（査読待ち初稿）} ISO26000は社会的責任に関するガイダンスであり、認証のための要求事項を規定した規格ではない。③がこの性格を誤っているため公式正答は③。ISOの一次資料でも認証用途ではないことを確認した。

\paragraph{他の選択肢（査読待ち初稿）} 他肢をすべて無条件に正しいと説明しない。①の政府除外の表現、④の原則の列挙範囲には一次資料との照合を要する疑義が残る。②⑤の歴史的な参加・連携の詳細も当時資料の確認を後工程に残す。

\par 参考：\cite{iso26000}。

\par 難易度 3（校準初稿）、優先度未評価。

\subsection{平成25年度 問6}
\label{q:h25-reg-apt-06}

\noindent\textbf{原典：} sources/H25/aptitude.pdf、PDF物理7頁。

\textbf{公式正答：} 5（single）。

\par\textit{原典本文の校正と独自解説は未作成。分類のみの登録項目。}

\subsection{令和元年度 問14}
\label{q:r01-reg-apt-14}

\noindent\textbf{原典：} sources/R01/aptitude.pdf、PDF物理14頁。

\textbf{公式正答：} 4（single）。

\par\textit{原典本文の校正と独自解説は未作成。分類のみの登録項目。}

\section{ユニバーサルデザイン}

\subsection{令和2年度 問11}
\label{q:r02-reg-apt-11}

\noindent\textbf{原典：} sources/R02/aptitude.pdf、PDF物理11頁。

\textbf{公式正答：} 1（single）。

\par\textit{原典本文の校正と独自解説は未作成。分類のみの登録項目。}

\section{未確定}

\subsection{平成24年度 問12}
\label{q:h24-reg-apt-12}

\noindent\textbf{原典：} sources/H24/aptitude.pdf、PDF物理10頁。

\textbf{公式正答：} 2（single）。

\par\textit{原典本文の校正と独自解説は未作成。分類のみの登録項目。}

\subsection{平成24年度 問2}
\label{q:h24-reg-apt-02}

\noindent\textbf{原典：} sources/H24/aptitude.pdf、PDF物理1頁。

\textbf{公式正答：} 3（single）。

\par\textit{原典本文の校正と独自解説は未作成。分類のみの登録項目。}

\subsection{平成25年度 問12}
\label{q:h25-reg-apt-12}

\noindent\textbf{原典：} sources/H25/aptitude.pdf、PDF物理13頁。

\textbf{公式正答：} 5（single）。

\par\textit{原典本文の校正と独自解説は未作成。分類のみの登録項目。}

\subsection{平成27年度 問15}
\label{q:h27-reg-apt-15}

\noindent\textbf{原典：} sources/H27/aptitude.pdf、PDF物理16頁。

\textbf{公式正答：} 4（single）。

\par\textit{原典本文の校正と独自解説は未作成。分類のみの登録項目。}

\subsection{平成27年度 問3}
\label{q:h27-reg-apt-03}

\noindent\textbf{原典：} sources/H27/aptitude.pdf、PDF物理3頁。

\textbf{公式正答：} 2（single）。

\par\textit{原典本文の校正と独自解説は未作成。分類のみの登録項目。}

\subsection{平成27年度 問6}
\label{q:h27-reg-apt-06}

\noindent\textbf{原典：} sources/H27/aptitude.pdf、PDF物理6頁。

\textbf{公式正答：} 3（single）。

\par\textit{原典本文の校正と独自解説は未作成。分類のみの登録項目。}

\subsection{平成29年度 問15}
\label{q:h29-reg-apt-15}

\noindent\textbf{原典：} sources/H29/aptitude.pdf、PDF物理14頁。

\textbf{公式正答：} 2（single）。

\par\textit{原典本文の校正と独自解説は未作成。分類のみの登録項目。}

\subsection{平成29年度 問3}
\label{q:h29-reg-apt-03}

\noindent\textbf{原典：} sources/H29/aptitude.pdf、PDF物理3頁。

\textbf{公式正答：} 3（single）。

\par\textit{原典本文の校正と独自解説は未作成。分類のみの登録項目。}

\subsection{平成30年度 問14}
\label{q:h30-reg-apt-14}

\noindent\textbf{原典：} sources/H30/aptitude.pdf、PDF物理14頁。

\textbf{公式正答：} ALL（all\_awarded）。

\par\textit{原典本文の校正と独自解説は未作成。分類のみの登録項目。}

\subsection{平成30年度 問15}
\label{q:h30-reg-apt-15}

\noindent\textbf{原典：} sources/H30/aptitude.pdf、PDF物理15頁。

\textbf{公式正答：} 4（single）。

\par\textit{原典本文の校正と独自解説は未作成。分類のみの登録項目。}

\subsection{平成30年度 問4}
\label{q:h30-reg-apt-04}

\noindent\textbf{原典：} sources/H30/aptitude.pdf、PDF物理4頁。

\textbf{公式正答：} 1（single）。

\par\textit{原典本文の校正と独自解説は未作成。分類のみの登録項目。}

\subsection{令和2年度 問2}
\label{q:r02-reg-apt-02}

\noindent\textbf{原典：} sources/R02/aptitude.pdf、PDF物理2頁。

\textbf{公式正答：} 2（single）。

\par\textit{原典本文の校正と独自解説は未作成。分類のみの登録項目。}

\subsection{令和5年度 問11}
\label{q:r05-reg-apt-11}

\noindent\textbf{原典：} sources/R05/aptitude.pdf、PDF物理11頁。

\textbf{公式正答：} 2（single）。

\par\textit{原典本文の校正と独自解説は未作成。分類のみの登録項目。}

\subsection{令和5年度 問2}
\label{q:r05-reg-apt-02}

\noindent\textbf{原典：} sources/R05/aptitude.pdf、PDF物理2頁。

\textbf{公式正答：} 3（single）。

\par\textit{原典本文の校正と独自解説は未作成。分類のみの登録項目。}

\subsection{令和5年度 問9}
\label{q:r05-reg-apt-09}

\noindent\textbf{原典：} sources/R05/aptitude.pdf、PDF物理9頁。

\textbf{公式正答：} 4（single）。

\par\textit{原典本文の校正と独自解説は未作成。分類のみの登録項目。}

\subsection{令和7年度 問3}
\label{q:r07-reg-apt-03}

\noindent\textbf{原典：} sources/R07/aptitude.pdf、PDF物理3頁。

\textbf{公式正答：} 5（single）。

\par\textit{原典本文の校正と独自解説は未作成。分類のみの登録項目。}
